% Source-derived chapter 01: Introduction
% Original source: very-stable-higgs-bundles-equivariant-multiplicity-mirror-symmetry
\chapter{Introduction}
\label{very-stable-higgs-bundles-equivariant-multiplicity-mirror-symmetry-chapter-01}
\source{very-stable-higgs-bundles-equivariant-multiplicity-mirror-symmetry}

\begin{remark}[Source section]
\label{very-stable-higgs-bundles-equivariant-multiplicity-mirror-symmetry-chapter-01-source}
\source{very-stable-higgs-bundles-equivariant-multiplicity-mirror-symmetry}
\source{very-stable-higgs-bundles-equivariant-multiplicity-mirror-symmetry}
This chapter reproduces the corresponding section of the retrieved
LaTeX source, including its definitions, statements, examples, and
proofs. It has not yet been translated into Lean.
\notready
\end{remark}

% The source section title was: Introduction
\stepcounter{subsection}
Drinfeld \cite{drinfeld}  and Laumon \cite{laumon} call a vector bundle on a smooth complex projective curve $C$ of genus $g$ a {\em very stable bundle} if there is no non-zero nilpotent Higgs field $\Phi\in H^0(C;\End(E)\otimes K)$ on $E$. Laumon \cite[Proposition 3.5]{laumon} proved that a very stable bundle
is semi-stable for $g>0$ and stable for $g>1$ in the sense of Mumford and that very stable bundles form a non-empty open subset in the moduli space of stable bundles on $C$. % For rank $2$ Pal and Pauly \cite{pal-pauly} showed, as was claimed by Drinfeld \cite{drinfeld}, that the complement of the locus of very stable bundles has codimension $1$. 

We will assume $g>1$, $d\in \Z$ and $n\in \Z_{>1}$. Let $\M:=\M^d_n$ denote the moduli space of semi-stable rank $n$ degree $d$ Higgs bundles on $C$ i.e. pairs $(E,\Phi)$ with $\Phi\in H^0(C;\End(E)\otimes K)$. It was constructed in Hitchin \cite{hitchin1} by gauge theoretic and in \cite{ nitsure,simpson} by algebraic geometric techniques. Pauly and P\'eon-Nieto \cite{pauly-peon-nieto} proved that $E$ is very stable if and only if \beq \label{stableup}U_{E}=\{(E,\Phi)|\Phi\in H^0(C;\End(E)\otimes K)\} \subset \M\eeq is closed in $\M$. 
\source{very-stable-higgs-bundles-equivariant-multiplicity-mirror-symmetry}

There is a generalization of $U_{E}\subset \M$ for more general Higgs bundles. Recall that there is a $\T$-action on $\M$ by $(E,\Phi)\mapsto (E,\lambda\Phi)$ for $\lambda\in \T$. For a stable Higgs bundle
$\calE=(E,\Phi)\in \M^{s\T}$ fixed by the $\T$-action we define its {\em upward flow} as \beq\label{upwardflow}W^+_\calE:=\{(E^\prime,\Phi^\prime)\in \M | \lim_{\lambda \to 0} (E^\prime,\lambda\Phi^\prime)=(E,\Phi)\}.\eeq
\source{very-stable-higgs-bundles-equivariant-multiplicity-mirror-symmetry}
In particular, for a rank $n$ degree $d$ vector bundle $E$ we have $W^+_{(E,0)}=U_E$ from \eqref{stableup}. 

We define a $\T$-fixed stable Higgs bundle $\calE\in \M^{s\T}$ to be a {\em very stable Higgs bundle} if the only Higgs bundle $(E^\prime,\Phi^\prime)\in W^+_\calE$ with nilpotent Higgs field is $\calE$.  Our first result is a generalization of \cite[Theorem 1.1]{pauly-peon-nieto}. \begin{theorem}\label{mainclosed}  $\calE\in \M^{s\T}$ is very stable if and only if $W^+_{\calE}\subset \M$ is closed. 
\source{very-stable-higgs-bundles-equivariant-multiplicity-mirror-symmetry}
\end{theorem}

From Laumon's result we know that a generic rank $n$ vector bundle is very stable, thus there always exists a very stable Higgs bundle of the form  $(E,0)$ with $E$ very stable, that is,  very stable Higgs  bundles in  the subvariety $\calN\in \pi_0(\M^\T)$ parametrizing semi-stable vector bundles on $C$. 

There are however many other fixed point components. Namely, when  $E=E_1\oplus\dots \oplus E_j$ is a direct sum of vector bundles and $\Phi|_{E_i}\to E_{i+1}K\subset EK$. We call $$(\rank(E_1),\rank(E_2),\dots,\rank(E_j))$$ -- an ordered partition of $n$ -- the {\em type} of the $\T$-invariant semistable Higgs bundle $(E,\Phi)$. For example, for a rank $n$ semistable bundle $E$ the semistable Higgs bundle $(E,0)\in \M^\T$ is of type $(n)$. 

In this paper we  classify all type $(1,\dots,1)$ very stable Higgs bundles:
\begin{theorem}
\source{very-stable-higgs-bundles-equivariant-multiplicity-mirror-symmetry}
\notready \label{maintype111} Let $$E=L_0\oplus \dots \oplus L_{n-1}$$ be a direct sum of line bundles $L_i$ on $C$. Furthermore, let the Higgs field be given by $$\Phi|_{L_{i-1}}=b_i : L_{i-1}\to L_{i}K\subset EK$$ where $$b_i\in \Hom(L_{i-1},L_{i}K)=H^0(C;L_{i-1}^*L_{i}K).$$ Assume that $(E,\Phi)$ is stable. Then $(E,\Phi)$ is very stable if and only if $$b=b_{n-1}\circ \dots \circ b_1\in H^0(C;L_0^*L_{n-1}K^{n-1})$$ has no repeated zero. 
\source{very-stable-higgs-bundles-equivariant-multiplicity-mirror-symmetry}
\end{theorem}

The proof of this proceeds by first verifying it in the case when $b$ has no zeroes. The corresponding Higgs bundle is called a uniformising Higgs bundle, the upward flow from which is nothing but a Hitchin section \cite{hitchin-lie}, which was already observed to be closed  in \cite{hitchin-lie}. In particular when $L_0\cong \calO_C$ then we get the {\em canonical uniformising Higgs bundle} \beq\label{canuni} \calE_0:=(\calO_C\oplus K^{-1}\oplus \dots \oplus K^{-n+1},\Phi|_{K^{-i+1}}=1:K^{-i+1}\to K^{-i}\otimes K).\eeq Then we apply Hecke transformations  at points of the curve to these Hitchin sections, and observe that we can reach all type $(1,\dots,1)$ upward flows with subsequent Hecke transformations, and can infer Theorem~\ref{maintype111} in turn. This argument is detailed in subsection~\ref{heckeverystable}.
\source{very-stable-higgs-bundles-equivariant-multiplicity-mirror-symmetry}

From Theorem~\ref{maintype111} we  deduce that in every type $(1,\dots,1)$ component $F\in\pi_0(\M^\T)$ there exists a very stable Higgs bundle. Thus  when $n=2$ every component, being either of type $(2)$ or type $(1,1)$, contains very stable Higgs bundles. However, already for rank $n=3$ there are type $(1,2)$ components $F\in\pi_0(\M^\T)$ which contain no very stable Higgs bundles. This can be seen by computing the multiplicity of a very stable component of the nilpotent cone, which we now explain. 

We recall the Hitchin map \cite{hitchin-stable} $$\begin{array}{cccc} h:&\M&\to &\calA=\oplus_{i=1}^n  H^0(C;K^i)\\ &(E,\Phi)&\mapsto & \det(\Phi-x)\end{array} $$ given by the characteristic polynomial. It is a proper completely integrable system \cite{nitsure,hitchin-stable} with respect to a natural symplectic structure on $\M$. In particular, $\dim(\calA)=\dim(\M)/2$ and the fibres of $h$ are Lagrangian subvarieties. The generic fibre is isomorphic to an Abelian variety -- the Jacobian of a  spectral curve -- while the most singular fibre $N:=h^{-1}(0)$, the so-called {\em nilpotent cone}, is non-reduced as a subscheme. Its components can be understood from the {\em core} of $\M$. Namely, for  $F\in {\pi_0(\M^\T)}$  one defines the {\em downward flow} $$ W^-_F:=\{(E^\prime,\Phi^\prime)\in \M | \lim_{\lambda \to \infty} (E^\prime,\lambda\Phi^\prime)\in F\}.$$ Note that the Hitchin map is proper and $\T$-equivariant with an appropriate weighted $\T$-action on $\calA$. It follows that the  core  $$\calC:=\bigcup_{F\in \pi_0(\M^\T)}W^-_F\cong N_{red}$$ is a proper subvariety isomorphic to $N_{red}$ the reduction of the nilpotent cone. Thus each component $N_F$ of the nilpotent cone $N$ is labelled by $F\in \pi_0(\M^{\T})$ and its reduction is isomorphic to ${N_F}_{red}\cong\overline{W^-_F}$. We denote by $m_F:=\ell(\calO_{N,N_F})$ the multiplicity of the component $N_F\subset N$. We have the following

\begin{theorem}
\source{very-stable-higgs-bundles-equivariant-multiplicity-mirror-symmetry}
\notready \label{mainmulti} Let $\calE\in \M^{s\T}$ be a very stable Higgs bundle. Denote by $T^+_\calE<T_\calE\M$ the positive part of the $\T$-module $T_\calE\M$ and by $F_\calE$ the $\T$-fixed point component containing $\calE$. We have  \begin{enumerate} \item the pushforward $h_*(\calO_{W^+_\calE})$ is a locally free sheaf on $\calA$ \item the multiplicity of the corresponding component $N_{F_\calE}$ of the nilpotent cone equals \beq \label{multi} m_{F_\calE}=\rank(h_*(\calO_{W^+_\calE}))=\left. { \chi_\T(\Sym(T^{+*}_\calE))\over \chi_\T(\Sym(\calA^*))}\right|_{t=1}\eeq  \item for generic $a\in \calA$  the intersection $W^+_\calE\cap h^{-1}(a)$ is transversal and has cardinality $m_{F_\calE}$. \end{enumerate}
\source{very-stable-higgs-bundles-equivariant-multiplicity-mirror-symmetry}
\end{theorem}
Here for a positive $\T$-module $V$ with weight space decomposition $V=\oplus_{\lambda>0} V_\lambda$ we have denoted the  $\T$-character of its symmetric algebra by 
\bes\chi_\T({\Sym(V^*)})= \sum_{\lambda\leq 0} \dim(\Sym(V^*)_\lambda)t^{-\lambda}= \prod_{\lambda>0}\frac{1}{( 1- t^{\lambda})^{\dim(V_\lambda)}} \in \Z[[t]],\ees where $t$ denotes the $\T$-character of weight $-1$. 
The computation of \eqref{multi} is straightforward and when it is non-integer valued we  deduce the non-existence of very stable Higgs bundles in a given component. This is the case for certain types  of $(1,2)$ components in  rank $3$, see Remark~\ref{counter}. However the integrity of $m_{F_\calE}$ is not a complete obstruction for existence of very stable bundles as examples in Remark~\ref{obstruction} for type $(1,3)$ in rank $4$ show.

It is straightforward to compute \eqref{multi} for a very stable Higgs bundle. Consequently, we  deduce the multiplicity of type $(n)$ and type $(1,\dots,1)$ components of the nilpotent cone. 

\begin{corollary}
\source{very-stable-higgs-bundles-equivariant-multiplicity-mirror-symmetry}
\notready \label{maincor} The multiplicity of $\calN$, the moduli space of semi-stable rank $n$ degree $d$ bundles, in the nilpotent cone is  $$m_\calN=2^{3g-3}3^{5g-5}\dots n^{(2n-1)(g-1)}.$$ Moreover let $N_{F_\calE}$ be the component of the nilpotent cone for a type $(1,\dots,1)$ Higgs bundle $\calE\in \M^{s\T}$ as in Theorem~\ref{maintype111}. Let $m_i=\deg(b_i)$. Then its multiplicity in the nilpotent cone is $$m_{F_{\calE}}=\prod_{i=1}^{n-1}{n\choose n-i}^{m_i}.$$
\source{very-stable-higgs-bundles-equivariant-multiplicity-mirror-symmetry}
\end{corollary}

The main motivation of this paper was to try to understand the mirror of upward flows  \eqref{upwardflow} in the framework of Kapustin and Witten \cite{kapustin-witten}. Fixing $\M:=\M_{n}^{-n(n-1)(g-1)}$ we expect that the mirror of the structure sheaf $\calO_{W^+_\calE}$ of the  Lagrangian upward flow ${W^+_\calE}\subset \M$ will be a hyperholomorphic vector bundle %\footnote{possibly twisted by a gerbe when $d\nmid n$} 
on $\M$ if and only if $\calE$ is very stable, in which case we expect the mirror to be of rank   $m_{F_\calE}$.  

More precisely,  for any $\calE\in \M^{s\T}$ we introduce the {\em virtual multiplicity} 
\beq\label{virmul}m_\calE(t):={ \chi_\T(\Sym(T^{+*}_\calE))\over \chi_\T(\Sym(\calA^*))}\eeq which is a priori  a rational function in $t$. In the course of proving Theorem~\ref{mainmulti} we find that for a very stable $\calE$ the pushforward $h_*(\calO_{W^+_\calE})$ is a $\T$-equivariant vector bundle on $\calA$, with  $\T$-equivariant fibre at $0$ satisfying $\chi_\T(h_*(\calO_{W^+_\calE})_0)=m_\calE(t)$. In particular, for a very stable $\calE$ the quantity $m_\calE(t)$ is a polynomial in $t$, which we call the {\em  equivariant multiplicity}. 
\source{very-stable-higgs-bundles-equivariant-multiplicity-mirror-symmetry}

Its significance for mirror symmetry is that we expect (see Remark~\ref{mirrorate0}) that  the mirror, denoted $\Lambda_\calE$, of a very stable upward flow $\calO_{W^+_\calE}$ should be a $\T$-equivariant vector bundle on $\M$ whose restriction to the canonical Hitchin section is isomorphic to $h_*(\calO_{W^+_\calE})$ and therefore its fibre at the canonical uniformising Higgs bundle $\calE_0$ should have $\T$-character \beq \label{mirrorexp}\chi_\T(\Lambda_\calE|_{\calE_0})=m_\calE(t).\eeq In particular, its rank should be $$\rank(\Lambda_\calE)=\dim(\Lambda_\calE|_{\calE_0})=m_\calE(1)=m_{F_\calE}$$ the multiplicity of $N_{F_\calE}$ in the nilpotent cone $N$. 
\source{very-stable-higgs-bundles-equivariant-multiplicity-mirror-symmetry}

In the case of type $(1,\dots,1)$ fixed points, where we have a complete classification of very stable Higgs bundles in Theorem~\ref{maintype111}, we  find a candidate for the mirror, which is supported by the following results sketched below. For details see Section~\ref{mirror}. 

Let $\delta=(\delta_0,\delta_1,\dots,\delta_{n-1})$ with $\delta_0:=c_{01}+\dots+c_{0m_0}$ a divisor of degree $\ell$, where we assume that either $c_{0j}\in C$ or $-c_{0j}\in C$.  For $i>0$ we have the effective divisor $\delta_i={c_{i1}+\dots+c_{im_{i}}}\in C^{[m_i]}$  of degree $m_i\geq 0$, where $c_{ij}\in C$ are points. This data gives rise to a type $(1,\dots,1)$ Higgs bundle  as in Theorem~\ref{maintype111} by setting $$L_i:=(\delta_0+\delta_1+\dots + \delta_{i})K^{-i}$$ and $$b_i:=s_{\delta_i}\in H^0(C;L_{i-1}^*L_{i}K)\cong H^0(C;\calO_C(\delta_i))$$ the defining section of $\calO_C(\delta_i)$. We  denote this Higgs bundle by  $\calE_\delta$. Let us assume that  $\calE_\delta\in \M^{s\T}$ is very stable, i.e. it is stable and the divisor $\delta_1+\dots + \delta_{n-1}$ is reduced. Then $W^+_{\calE_\delta}\subset \M$ is a closed Lagrangian, and we denote its structure sheaf $$\calL_\delta:=\calO_{W^+_{\calE_\delta}}.$$

We construct  \beq\label{lambdaintro}\Lambda_\delta:= \bigotimes_{i=0}^{n-1}\bigotimes_{j=0}^{m_i} \Lambda^{n-i}(\bE_{c_{ij}}),\eeq  using a certain (twisted by a gerbe and appropriately normalised) universal Higgs bundle $(\bE,\bPhi)$ on $C\times \M^s$ and denote $\bE_{c_{ij}}=\bE|_{\M^s\times{\{c_{ij}\}}}$ when $c_{ij}\in C$ and $\bE_{c_{0j}}=\bE^*|_{\M^s\times{\{-c_{0j}\}}}$, when $-c_{0j}\in C$. When $n$ divides $\sum_{i=1}^{n-1} (n-i)m_i$ \eqref{lambdaintro} is a well-defined vector bundle of rank $m_{\calE_\delta}$.
\source{very-stable-higgs-bundles-equivariant-multiplicity-mirror-symmetry}

For simplicity we assume that $$d=n\ell+\sum_{i=1}^{n-1}(n-i)m_i-n(n-1)(g-1)=-n(n-1)(g-1).$$ In this case for a generic $a\in \calA^\#\subset \calA$  we can identify $h^{-1}(a)=J(C_a)$ with the Jacobian of the smooth spectral curve $C_a\subset T^*C$. Define $$\begin{array}{cccc}\iota:&\M&\to& \M \\ &(E,\Phi)&\mapsto& (E^*K^{1-n},\Phi^T)\end{array},$$ which on $h^{-1}(a)$ will induce the inverse map. In particular, we will have $\iota(\calE_\delta)=\calE_{\delta^\iota}$ where $\delta^\iota=(\delta^\iota_0,\delta^\iota_1,\dots,\delta^\iota_{n-1})=(-\delta_{n-1}-\dots-\delta_0, \delta_{n-1},\dots,\delta_1).$

For $\M^\#:=h^{-1}(\calA^\#)$ we  construct a relative Poincar\'e bundle $\tilde{\bP}$ on the fibre product $\M^\#\times_{\calA^\#}\M^\#$, which has fibre $J(C_a)\times J(C_a)$ over $a\in \calA^\#$. 
We then have the following 
\begin{theorem}
\source{very-stable-higgs-bundles-equivariant-multiplicity-mirror-symmetry}
\notready \label{mainmirror} Let $\calE_\delta\in \M^{s\T}$ be a very stable Higgs bundle of type $(1,\dots,1)$ and degree $-n(n-1)(g-1)$.  The relative Fourier-Mukai transform over $\calA^\#$ satisfies $$S(\calL_\delta|_{\M^\#})=(\pi_2)_*(\pi_1^*(\calL_\delta|_{\M^\#})\otimes \tilde{\bP})= \Lambda_\delta|_{\M^\#}$$
\source{very-stable-higgs-bundles-equivariant-multiplicity-mirror-symmetry}
	
	and 	$$S(\Lambda_\delta|_{\M^\#})=\calL_{\delta^\iota}[-n^2(g-1)-1]|_{\M^\#}$$
	
	Moreover,   $\Lambda_\delta$ extends to  a  vector bundle with a hyperholomorphic connection on $\M^s$. 
\end{theorem} 

This theorem gives the mirror of $\calL_\delta$  as a hyperholomorphic vector bundle of rank $m_{F_{\calE_\delta}}$. The following will give evidence of the mirror relationship globally, in particular over the nilpotent cone.  To motivate it recall that mirror symmetry should be an equivalence of categories. Not just objects should match but morphisms between them. If two objects have a $\T$-equivariant structure so will the vector space of morphisms between them. Thus its $\T$-character should agree with the $\T$-character of the vector space of morphisms between the mirror objects.

To formalise this for two $\T$-equivariant coherent sheaves $\calF_1$ and $\calF_2$ on $\M$  we denote the {\em equivariant Euler pairing} as

$$\chi_\T(\M;\calF_1,\calF_2)=\sum_{k,l} \dim(H^k({ R} {H\!om}(\calF_1,\calF_2))_l) (-1)^k t^{-l},$$ which
on a semi-projective variety like $\M$ we expect to be a Laurent series in $\Z((t))$. 

\begin{theorem}
\source{very-stable-higgs-bundles-equivariant-multiplicity-mirror-symmetry}
\notready \label{maineuler}  Let $\calE_\delta,\calE_{\delta^\prime}\in \M^{s\T}$ be two type $(1,\dots,1)$ very stable Higgs bundles of degree $-n(n-1)(g-1)$. Then there exist $\T$-equivariant structures on all coherent  sheaves of the form $\calL_\delta$ and $\Lambda_\delta$  such that $$\chi_\T\left(\M;\calL_{\delta^\prime},\Lambda_{\delta}\right)=\chi_\T(\det(\C_{1}\otimes\calA^*)) \chi_\T\left(\M;\Lambda_{\delta^\prime},\calL_{\delta^\iota}[-n^2(g-1)-1]
\source{very-stable-higgs-bundles-equivariant-multiplicity-mirror-symmetry}
	\right). $$
	
\end{theorem}

Thus up to the $q$-power $\chi_\T(\det(\C_{1}\otimes\calA^*))$, where $\C_1$ denotes the weight $1$ $\T$-character,  we find that the $\T$-character of morphisms between our Lagrangian and hyperholomorphic branes agree with that of the conjectured mirror branes. 

The content of  the paper is as follows. In Section~\ref{bb} we work out from scratch the Bialynicki-Birula theory for a semi-projective variety,  to conclude that the upward flows and the core of a semiprojective variety with a homogeneity $1$ symplectic form are Lagrangian. Here we  introduce very stable upward flows and prove some  properties including the general form of Theorem~\ref{mainclosed}.  In Section~\ref{bbhiggs} we determine which Higgs bundles belong to which upward and downward flows in the moduli space of Higgs bundles. In Section~\ref{verystable} we discuss very stable Higgs bundles and prove part $1$ of Theorem~\ref{mainmulti}. Then we examine in detail the Hecke transformation  of a  type $(1,\dots,1)$ $\T$-fixed Higgs bundle and apply this to classify all very stable Higgs bundles of this type proving Theorem~\ref{maintype111}. Moreover we show that the Hecke transform takes a Lagrangian subvariety to another in general. In Section~\ref{multiplicities} we investigate the multiplicity of a very stable component of the nilpotent cone, by deriving an explicit linear algebra formula for it, and study the implications. Here we prove the rest of Theorem~\ref{mainmulti} and Corollary~\ref{maincor}.  In Section~\ref{mirror} we study the relative Fourier-Mukai transform of our type $(1,\dots,1)$ very stable upward flows proving the first part of  Theorem~\ref{mainmirror}, and prove an agreement of $\T$-equivariant Euler forms for pairs of mirror branes settling Theorem~\ref{maineuler}. In Section~\ref{hyper} using the twistor space construction of the moduli space of Higgs bundles we endow the universal bundle at a point in $C$ with a hyperholomorphic connection and from the expression (\ref{lambdaintro}) prove  the second part of Theorem~\ref{mainmirror}. 

In the final Section~\ref{further} we discuss three issues relevant to further progress. First in Subsection~\ref{simple} we compute the virtual multiplicity \eqref{virmul} for the analogue of  the type $(1,\dots,1)$ fixed points for a simple group $G$ instead of $\GL_n$. This is the case where the $\T$-action at the fixed point set has the same form as the Hitchin section for the group $G$ as in \cite{hitchin-lie}, which is of course a closed upward flow. The $b_i$ which we studied for the general linear group are now replaced by sections $b_1,\dots, b_{\ell}$ of line bundles associated to the simple roots of $G$.
We find that unlike  the $\GL_n$ case considered in the body of this paper, the virtual multiplicity is rarely a polynomial. We prove in Proposition~\ref{cominint} that for {\em cominuscule} Higgs bundles where $b_i$ are nonvanishing sections except at  cominuscule roots it is in fact a polynomial.  In this case we offer Conjecture~\ref{cominmirror} that under precise conditions these are very stable and find a potential mirror in terms of the minuscule representations of the Langlands dual group, satisfying the expectation of \eqref{mirrorexp}. Secondly in Subsection~\ref{multiple} we study the first non-trivial wobbly type $(1,1)$ $\SL_2$ Higgs bundle where $b$ has a single double $0$. We find the closure of its upward flow and formulate a conjecture about the mirror of the universal  $\SO_3$ bundle  in the adjoint representation. Finally in Subsection~\ref{cotfib} we discuss the problem of the mirror of a very stable cotangent fibre to the moduli of semi-stable vector bundles. We relate this to considerations of Donagi--Pantev and to Drinfeld--Laumon in the geometric Langlands correspondence. 



%From the existence of very stable bundles \cite[Proposition 3.5]{laumon}, i.e. very stable Higgs bundles in $\calN\in  \M^\T$ we conclude - see Remark~\ref{multN} - that the multiplicity of the moduli space of semi-stable bundles $\calN\subset N\subset \M$ in the nilpotent cone is \beq\label{multn}m_\calN=2^{3g-3}3^{5g-5}\cdots n^{(2n-1)(g-1)}.\eeq
%In fact, this number \eqref{multn} has appeared in \cite{beauville-etal} as the degree of the map \beq \label{bnrdom}\pi_a: h^{-1}(a)\cap T^*\calN \to \calN\eeq for a generic $a$, which  agrees with the cardinality of a transversal intersection $U_{E}\cap h^{-1}(a)$ for a very stable bundle $E$, which in turn agrees with $m_{\calN}$ by Theorem~\ref{mainmulti}.2. 

%The number \eqref{multn} also appears in \cite[Conjecture 6.2.1]{laumon2} as the conjectured rank of cuspidal automorphic sheaves on the moduli space $\calN^{vs}\subset \calN^{s}\subset \calN$ of very stable vector bundles in the Geometric Langlands correspondence for function fields. This was in fact the main motivation for Drinfeld \cite{drinfeld} to introduce very stable bundles to analyze cuspidal automorphic sheaves, he constructed in \cite{drinfeld1} for $\GL_2$. Laumon after \cite[Conjecture 6.2.1]{laumon2} speculated that there should be a bijection between  rank $n$ complex local systems on the complex curve $C$ and  rank $m_\calN$ local systems on $\calN^{vs}$ whose perverse extionsion to $\calN$ (in fact to the whole stack of rank $n$ vector bundles) has its characteristic variety contained  in the nilpotent cone.  

%We can see the relevance of very stable bundles and the multiplicity in \eqref{multn} in the conjectural mirror symmetry approach to the Geometric Langlands correspondence \cite{hausel0,donagi-pantev}. Here one starts with an irreducible rank $n$ complex local system on $C$. By the non-abelian Hodge theorem of \cite{hitchin,simpson} this is identified with a rank $n$ stable Higgs bundle $\calE\in \M^s$. Assuming that $h(\calE)=a\in \calA$ is generic we can compute the Fourier-Mukai transform $S(\calO_{\calE})$ of the structure sheaf of the point $\calE\in h^{-1}(a)$ yielding a degree $0$ line bundle on  $h^{-1}(0)\cong J(C_a)$  the Jacobian of a spectral curve $C_a$. In turn, one considers the map $\pi_a$ of \eqref{bnrdom} generically as a spectral curve for $\calN$, and pushing forward $S(\calO_\calE)$ by $\pi_a$ one gets a quasi-coherent Higgs sheaf
%on $\calN$, which will be a Higgs bundle  exactly on $\calN^{vs}$ of rank equalling the degree of $\pi_a$ i.e. $m_\calN$ in \eqref{multn}.  To conclude, by an extension of Simpson's non-abelian Hodge theorem, due to Mochizuki \cite{mochizuki}, should give the required $D$-module on $\calN$. 

%In this mirror symmetry approach we are also interested in the mirror of more general coherent sheaves
%on $\M$ than the skyscraper sheaves above. In the physics motivated view point of  \cite{kapustin-witten} the structure sheaves of Lagrangian subvarieties in $\M$ should be considered so-called $B,A,A$ branes, which should be mirror to $B,B,B$ branes which live on  vector bundles on hyperk\"ahler subvarieties of $\M$. One example is the  mirror pair above: the $B,B,B$ brane of the skypscraper sheaf $\calO_\calE$ mirror to the $B,A,A$ brane given by a line bundle on $h^{-1}(a)$. 

%Another immediate example is the conjectured mirror pair of a $B,B,B$ brane of the structure sheaf $\calO_\M$ with the $B,A,A$ brane of the identity section $U_0$ of the Hitchin map. We notice that the identity section of the Hitchin map $U_0=U_{\calE_0}$ coincides with the upward flow from the uniformizing Higgs bundle $\calE_0$. 
%More generally we expect that the mirrors of the Lagrangian very stable upward flows $\calO_{U_\calE}$ will be vector bundles on the total space $\M$. Our next result is the

\vskip.5cm




{\noindent \bf Acknowledgements.} We would like to thank Brian Collier, Davide Gaiotto, Peter Gothen, Jochen Heinloth, Daniel Huybrechts, Quoc Ho, Joel Kamnitzer, G\'erard Laumon, Luca Migliorini, Alexander Minets, Brent Pym,  Peng Shan, Carlos Simpson, Andr\'as Szenes, Fernando R. Villegas, Richard Wentworth, Edward Witten and K\={o}ta Yoshioka for interesting comments and discussions. Most of all we are grateful for a long list of very helpful comments by the referee. We would also like to thank the organizers of the Summer School on Higgs bundles in Hamburg in September 2018, where the authors and Richard Wentworth were  giving lectures and where the work in this paper started by considering the mirror of the Lagrangian upward flows $W^+_\calE$ investigated in \cite{collier-wentworth}. The second author wishes to thank EPSRC and ICMAT for support.
